% ============================================================
% EXAMEN DIAGNÓSTICO: TERMODINÁMICA
% Curso de Oriente - CCH Oriente
% Enero 2026
% ============================================================
% INSTRUCCIONES: Pegar este contenido en tu plantilla de Overleaf
% después de la definición de entornos y antes de \end{document}
% ============================================================
\renewcommand{\tema}{Termodinámica}
% ============================================================

\twocolumn

% ============================================================
\section*{EXAMEN DIAGNÓSTICO — TERMODINÁMICA}
% ============================================================

% ============================================================
% PREGUNTAS
% ============================================================

\begin{preguntabox}
\subsection*{Pregunta 1}
Un estudiante sostiene un cubo de hielo en la mano y afirma: ``el frío del hielo pasa a mi mano''. ¿Qué error conceptual comete?

\begin{enumerate}
    \item[A)] Ninguno; el frío sí se transfiere de cuerpos fríos a calientes.
    \item[B)] El calor fluye de su mano hacia el hielo, no al revés.
    \item[C)] La temperatura del hielo es igual a la de su mano.
    \item[D)] El calor y la temperatura son la misma magnitud.
\end{enumerate}
\end{preguntabox}

\begin{preguntabox}
\subsection*{Pregunta 2}
Dos objetos se ponen en contacto térmico y, tras un tiempo, dejan de intercambiar calor. ¿Qué ley de la termodinámica describe esta situación?

\begin{enumerate}
    \item[A)] Primera ley de la termodinámica.
    \item[B)] Segunda ley de la termodinámica.
    \item[C)] Ley cero de la termodinámica.
    \item[D)] Ley de Gay-Lussac.
\end{enumerate}
\end{preguntabox}

\begin{preguntabox}
\subsection*{Pregunta 3}
¿Cuál de los siguientes mecanismos de transferencia de calor \textbf{no} requiere un medio material para propagarse?

\begin{enumerate}
    \item[A)] Conducción.
    \item[B)] Convección.
    \item[C)] Radiación.
    \item[D)] Conducción y convección.
\end{enumerate}
\end{preguntabox}

\begin{preguntabox}
\subsection*{Pregunta 4}
La temperatura corporal normal es 37$^\circ$C. ¿Cuál es su equivalente en la escala Kelvin?

\begin{enumerate}
    \item[A)] 236 K
    \item[B)] 273 K
    \item[C)] 310 K
    \item[D)] 337 K
\end{enumerate}
\end{preguntabox}

\begin{preguntabox}
\subsection*{Pregunta 5}
Se calientan 2 kg de agua desde 20$^\circ$C hasta 35$^\circ$C. Sabiendo que el calor específico del agua es 1 cal/(g$\cdot^\circ$C), ¿cuánto calor se necesita?

\begin{enumerate}
    \item[A)] 15 kcal
    \item[B)] 30 kcal
    \item[C)] 70 kcal
    \item[D)] 150 kcal
\end{enumerate}
\end{preguntabox}

\begin{preguntabox}
\subsection*{Pregunta 6}
La primera ley de la termodinámica establece que la variación de energía interna de un sistema es igual a:

\begin{enumerate}
    \item[A)] El calor absorbido más el trabajo realizado sobre el sistema.
    \item[B)] El calor absorbido menos el trabajo realizado por el sistema.
    \item[C)] Solo el trabajo realizado por el sistema.
    \item[D)] Solo el calor cedido al ambiente.
\end{enumerate}
\end{preguntabox}

\begin{preguntabox}
\subsection*{Pregunta 7}
Un gas ocupa 3 L a una presión de 2 atm. Si la presión disminuye a 1 atm manteniendo la temperatura constante, ¿cuál es el nuevo volumen?

\begin{enumerate}
    \item[A)] 1.5 L
    \item[B)] 3 L
    \item[C)] 6 L
    \item[D)] 9 L
\end{enumerate}
\end{preguntabox}

\begin{preguntabox}
\subsection*{Pregunta 8}
Un gas ideal ocupa 4 L a 300 K. Si se calienta a presión constante hasta 450 K, ¿qué volumen ocupa?

\begin{enumerate}
    \item[A)] 2 L
    \item[B)] 4 L
    \item[C)] 6 L
    \item[D)] 8 L
\end{enumerate}
\end{preguntabox}

\begin{preguntabox}
\subsection*{Pregunta 9}
Un recipiente rígido contiene gas a 3 atm y 300 K. Al calentarlo hasta 400 K, ¿cuál es la nueva presión?

\begin{enumerate}
    \item[A)] 2.25 atm
    \item[B)] 3 atm
    \item[C)] 4 atm
    \item[D)] 5 atm
\end{enumerate}
\end{preguntabox}

\begin{preguntabox}
\subsection*{Pregunta 10}
Una pieza de aluminio de 500 g se calienta de 20$^\circ$C a 100$^\circ$C. El calor específico del aluminio es 0.5 cal/(g$\cdot^\circ$C). ¿Cuánto calor absorbe?

\begin{enumerate}
    \item[A)] 10 kcal
    \item[B)] 20 kcal
    \item[C)] 40 kcal
    \item[D)] 25 kcal
\end{enumerate}
\end{preguntabox}

% ============================================================
\newpage
\onecolumn
\section*{RESPUESTAS Y RETROALIMENTACIÓN}
% ============================================================

\begin{respuestabox}
\subsection*{Pregunta 1}
\textcolor{verdecorrecto}{\textbf{Respuesta correcta: B) El calor fluye de su mano hacia el hielo, no al revés.}}

\textbf{Justificación:}\\
El calor es energía en tránsito que fluye \textbf{siempre} del cuerpo de mayor temperatura al de menor temperatura. La mano está a $\approx$37$^\circ$C y el hielo a 0$^\circ$C, por lo que el calor va de la mano al hielo. La sensación de ``frío'' es la percepción de que la mano pierde energía, no que el hielo ``envíe frío''.

\textbf{Respuestas incorrectas:}
\begin{itemize}
    \item \textbf{A) Ninguno:} El frío no es una magnitud que se transfiera; solo el calor se transfiere, y siempre de mayor a menor temperatura.
    \item \textbf{C) La temperatura es igual:} Si las temperaturas fueran iguales, no habría flujo de calor; el hielo se derretiría porque sí hay diferencia de temperatura.
    \item \textbf{D) Calor y temperatura son lo mismo:} El calor es energía en tránsito (J o cal); la temperatura es una propiedad del estado del sistema (K, $^\circ$C). Son magnitudes distintas.
\end{itemize}
\end{respuestabox}

\begin{respuestabox}
\subsection*{Pregunta 2}
\textcolor{verdecorrecto}{\textbf{Respuesta correcta: C) Ley cero de la termodinámica.}}

\textbf{Justificación:}\\
La \textbf{ley cero} establece que si dos sistemas están en equilibrio térmico con un tercero, están en equilibrio térmico entre sí. Cuando dos cuerpos en contacto dejan de intercambiar calor, han alcanzado la misma temperatura: se dice que están en \textbf{equilibrio térmico}. Este principio es la base conceptual de la medición de temperatura con termómetros.

\textbf{Respuestas incorrectas:}
\begin{itemize}
    \item \textbf{A) Primera ley:} Describe la conservación de la energía interna ($\Delta U = Q - W$), no el equilibrio térmico entre cuerpos.
    \item \textbf{B) Segunda ley:} Establece la dirección irreversible de los procesos y el concepto de entropía.
    \item \textbf{D) Gay-Lussac:} Es una ley de gases ideales que relaciona presión y temperatura a volumen constante; no tiene que ver con equilibrio térmico.
\end{itemize}
\end{respuestabox}

\begin{respuestabox}
\subsection*{Pregunta 3}
\textcolor{verdecorrecto}{\textbf{Respuesta correcta: C) Radiación.}}

\textbf{Justificación:}\\
La \textbf{radiación} es la transferencia de calor mediante ondas electromagnéticas (principalmente infrarrojas). Al ser ondas electromagnéticas, se propagan en el vacío sin necesitar un medio material. Por eso el calor del Sol llega a la Tierra atravesando el espacio. En cambio, la \textbf{conducción} requiere contacto directo entre partículas, y la \textbf{convección} requiere el movimiento de un fluido.

\textbf{Respuestas incorrectas:}
\begin{itemize}
    \item \textbf{A) Conducción:} Ocurre por vibración y choque entre moléculas adyacentes; necesita un medio sólido o fluido en contacto.
    \item \textbf{B) Convección:} Ocurre por el movimiento de masas de fluido (corrientes); requiere un fluido como medio.
    \item \textbf{D) Conducción y convección:} Ambas requieren un medio material, por lo que ninguna puede propagarse en el vacío.
\end{itemize}
\end{respuestabox}

\begin{respuestabox}
\subsection*{Pregunta 4}
\textcolor{verdecorrecto}{\textbf{Respuesta correcta: C) 310 K}}

\textbf{Justificación:}\\
La conversión de Celsius a Kelvin es:
$$T(\text{K}) = T(^\circ\text{C}) + 273$$
Sustituyendo:
$$T(\text{K}) = 37\,^\circ\text{C} + 273 = 310\,\text{K}$$

\textbf{Respuestas incorrectas:}
\begin{itemize}
    \item \textbf{A) 236 K:} Resulta de restar 273 en lugar de sumar: $37 - 273 = -236$, lo cual representa una temperatura por debajo del cero absoluto, que es imposible.
    \item \textbf{B) 273 K:} Es la conversión de 0$^\circ$C a Kelvin, no de 37$^\circ$C. El alumno confundió el punto de referencia.
    \item \textbf{D) 337 K:} Resulta de sumar 300 en lugar de 273: error por usar una constante incorrecta de conversión.
\end{itemize}
\end{respuestabox}

\begin{respuestabox}
\subsection*{Pregunta 5}
\textcolor{verdecorrecto}{\textbf{Respuesta correcta: B) 30 kcal}}

\textbf{Justificación:}\\
Se aplica la fórmula del calor:
$$Q = m \cdot c \cdot \Delta T$$
donde $m = 2\,\text{kg} = 2000\,\text{g}$, $c = 1\,\dfrac{\text{cal}}{\text{g}\cdot^\circ\text{C}}$ y $\Delta T = 35\,^\circ\text{C} - 20\,^\circ\text{C} = 15\,^\circ\text{C}$.

$$Q = 2000\,\text{g} \times 1\,\frac{\text{cal}}{\text{g}\cdot^\circ\text{C}} \times 15\,^\circ\text{C} = 30{,}000\,\text{cal} = 30\,\text{kcal}$$

\textbf{Respuestas incorrectas:}
\begin{itemize}
    \item \textbf{A) 15 kcal:} Olvidó convertir kg a gramos; usó $m = 2$ en lugar de $m = 2000\,\text{g}$, obteniendo $Q = 2 \times 1 \times 15 = 30\,\text{cal}$, y luego interpretó eso como 15 kcal dividiendo entre 2.
    \item \textbf{C) 70 kcal:} Usó $\Delta T = 35\,^\circ\text{C}$ en lugar de la diferencia $15\,^\circ\text{C}$.
    \item \textbf{D) 150 kcal:} Multiplicó la masa en kg directamente por la temperatura final (35) en lugar de por el $\Delta T$ correcto.
\end{itemize}
\end{respuestabox}

\begin{respuestabox}
\subsection*{Pregunta 6}
\textcolor{verdecorrecto}{\textbf{Respuesta correcta: B) El calor absorbido menos el trabajo realizado por el sistema.}}

\textbf{Justificación:}\\
La primera ley de la termodinámica expresa la conservación de la energía:
$$\Delta U = Q - W$$
donde $\Delta U$ es la variación de energía interna, $Q$ es el calor absorbido por el sistema y $W$ es el trabajo realizado \textbf{por} el sistema. Si el sistema absorbe calor ($Q > 0$) y realiza trabajo ($W > 0$), una parte del calor se convierte en trabajo y el resto incrementa la energía interna.

\textbf{Respuestas incorrectas:}
\begin{itemize}
    \item \textbf{A) Calor más trabajo:} Confunde el signo del trabajo. $\Delta U = Q + W$ sería correcto si $W$ fuera el trabajo realizado \textbf{sobre} el sistema, pero la convención estándar define $W$ como trabajo del sistema hacia el exterior.
    \item \textbf{C) Solo el trabajo:} Omite la contribución del calor, que es fundamental en la primera ley.
    \item \textbf{D) Solo el calor cedido:} El calor cedido implicaría $Q < 0$; además, la ley involucra tanto $Q$ como $W$.
\end{itemize}
\end{respuestabox}

\begin{respuestabox}
\subsection*{Pregunta 7}
\textcolor{verdecorrecto}{\textbf{Respuesta correcta: C) 6 L}}

\textbf{Justificación:}\\
A temperatura constante se aplica la \textbf{ley de Boyle}: $P_1 V_1 = P_2 V_2$.

Despejando $V_2$:
$$V_2 = \frac{P_1 \cdot V_1}{P_2} = \frac{2\,\text{atm} \times 3\,\text{L}}{1\,\text{atm}} = 6\,\text{L}$$

Al reducirse la presión a la mitad, el volumen se duplica.

\textbf{Respuestas incorrectas:}
\begin{itemize}
    \item \textbf{A) 1.5 L:} Invirtió la relación; dividió el volumen entre 2 cuando debía multiplicarlo.
    \item \textbf{B) 3 L:} Supuso que el volumen no cambia, ignorando la relación inversa entre presión y volumen.
    \item \textbf{D) 9 L:} Sumó los valores de presión y volumen en lugar de aplicar la proporcionalidad inversa.
\end{itemize}
\end{respuestabox}

\begin{respuestabox}
\subsection*{Pregunta 8}
\textcolor{verdecorrecto}{\textbf{Respuesta correcta: C) 6 L}}

\textbf{Justificación:}\\
A presión constante se aplica la \textbf{ley de Charles}: $\dfrac{V_1}{T_1} = \dfrac{V_2}{T_2}$.

Despejando $V_2$:
$$V_2 = V_1 \cdot \frac{T_2}{T_1} = 4\,\text{L} \times \frac{450\,\text{K}}{300\,\text{K}} = 4\,\text{L} \times 1.5 = 6\,\text{L}$$

\textbf{Importante:} Las temperaturas deben estar en Kelvin para aplicar esta ley.

\textbf{Respuestas incorrectas:}
\begin{itemize}
    \item \textbf{A) 2 L:} Invirtió la razón de temperaturas: usó $T_1/T_2$ en lugar de $T_2/T_1$, obteniendo $4 \times (300/450) \approx 2.67\,\text{L}$, y redondeó incorrectamente.
    \item \textbf{B) 4 L:} Supuso que el volumen no cambia al variar la temperatura, omitiendo la relación directa entre $V$ y $T$.
    \item \textbf{D) 8 L:} Usó la diferencia de temperaturas ($450 - 300 = 150\,\text{K}$) como factor multiplicativo, sin base física.
\end{itemize}
\end{respuestabox}

\begin{respuestabox}
\subsection*{Pregunta 9}
\textcolor{verdecorrecto}{\textbf{Respuesta correcta: C) 4 atm}}

\textbf{Justificación:}\\
En un recipiente rígido el volumen es constante, por lo que se aplica la \textbf{ley de Gay-Lussac}: $\dfrac{P_1}{T_1} = \dfrac{P_2}{T_2}$.

Despejando $P_2$:
$$P_2 = P_1 \cdot \frac{T_2}{T_1} = 3\,\text{atm} \times \frac{400\,\text{K}}{300\,\text{K}} = 3\,\text{atm} \times \frac{4}{3} = 4\,\text{atm}$$

\textbf{Respuestas incorrectas:}
\begin{itemize}
    \item \textbf{A) 2.25 atm:} Invirtió la razón de temperaturas: $3 \times (300/400) = 2.25\,\text{atm}$. Al calentar un gas a volumen fijo, la presión sube, no baja.
    \item \textbf{B) 3 atm:} Supuso que la presión no cambia, ignorando la relación directa entre $P$ y $T$ a volumen constante.
    \item \textbf{D) 5 atm:} Sumó la diferencia de temperaturas ($400 - 300 = 100$) a la presión original sin justificación física.
\end{itemize}
\end{respuestabox}

\begin{respuestabox}
\subsection*{Pregunta 10}
\textcolor{verdecorrecto}{\textbf{Respuesta correcta: B) 20 kcal}}

\textbf{Justificación:}\\
Se aplica $Q = m \cdot c \cdot \Delta T$:
\begin{itemize}
    \item $m = 500\,\text{g}$
    \item $c = 0.5\,\dfrac{\text{cal}}{\text{g}\cdot^\circ\text{C}}$
    \item $\Delta T = 100\,^\circ\text{C} - 20\,^\circ\text{C} = 80\,^\circ\text{C}$
\end{itemize}

$$Q = 500\,\text{g} \times 0.5\,\frac{\text{cal}}{\text{g}\cdot^\circ\text{C}} \times 80\,^\circ\text{C} = 20{,}000\,\text{cal} = 20\,\text{kcal}$$

\textbf{Respuestas incorrectas:}
\begin{itemize}
    \item \textbf{A) 10 kcal:} Usó la temperatura final (100$^\circ$C) como $\Delta T$ y dividió entre dos sin justificación; o bien olvidó multiplicar por $c = 0.5$ al usar solo la mitad.
    \item \textbf{C) 40 kcal:} Usó $c = 1\,\text{cal/(g·}^\circ\text{C)}$ (valor del agua) en lugar de $c = 0.5\,\text{cal/(g·}^\circ\text{C)}$ del aluminio.
    \item \textbf{D) 25 kcal:} Usó $\Delta T = 100\,^\circ\text{C}$ en lugar de la diferencia correcta de $80\,^\circ\text{C}$, obteniendo $500 \times 0.5 \times 100 = 25{,}000\,\text{cal}$.
\end{itemize}
\end{respuestabox}
